\printedpage{29}
% head: A. Grothendieck : Dérivateurs, Ch. XIII, Édition des Universités de Montpellier II et Paris VII
(We apply \texttt{[the]} corollary 1 (p. 47) to $\tilde{\alpha} : \boldsymbol{\Delta}_1 \longrightarrow I$.)

\textbf{Proposition 3.} (\footnote[28]{Compare with the lemma p. 63 and its proof (more exactly, the corollary of the lemma).}) \textit{Let $X$ be in $\mathcal{M}(\boldsymbol{\Delta}_1)$, $X = (X_0 \xleftarrow{\alpha} X_1)$. Then $X$ is in $\mathcal{M}(\boldsymbol{\Delta}_1)_{\mathrm{f}}$ if and only if $X_0$, $X_1$} \texttt{[are]} \textit{in $\mathcal{M}_{\mathrm{f}}$ and $\alpha \in F$.}

\textbf{Corollary 2.} \textit{Let $I' \overset{u}{\hookrightarrow} I$ be} \texttt{[an]} open immersion. \textit{Then for $X \in \operatorname{Ob}\mathcal{M}(I)_{\mathrm{f}}$, $\varprojlim\limits_{I} X(i) \longrightarrow \varprojlim\limits_{I'} X(i)$ is in $F$.}

\section*{5\quad Quillenizable categories of models}

I return to the general case of a category of models $(\mathcal{M},W)$.

% marginal vertical bar printed beside the following Définition 1
\textbf{Definition 1.} We say that $(\mathcal{M},W)$ is \textit{Quillenizable} if there exists an \textit{exact} weak Quillen structure $(C,TC,F,TF)$ for $\mathcal{M}$.

That is a self-dual condition, but one which I know how to verify in practice only by means of either left or right exactness -- thus by a stronger condition which destroys the symmetry. It is to preserve it as far as possible that I introduce the preceding definition, as well as the

\textbf{Definition 2.} Let, in addition to $(\mathcal{M},W)$, Diag $\subset$ Cat be a full subcategory of Cat. We say that
\msp{page 51}$W$ is \textit{Quillenizable on} Diag, if for every $I$ in Diag, $(\mathcal{M}(I),W(I))$ is Quillenizable.

In practice, I will use this definition in the case where Diag $=$ Cat, or else, that it satisfy at least the stability properties useful for the domain of a derivator: stable under finite $\varprojlim$, under finite sums -- possibly under finite $\varinjlim$, etc (\footnote[29]{I will assume that Diag contains at least the categories coming from finite ordered sets.}). % [sic] eventuellement -> éventuellement, missing accent

\textbf{Example 1:} That of paragraphs 3, 4, where we isolated conditions on $(\mathcal{M},W)$, with $\mathcal{M}$ a ``big'' category and $\mathcal{M}$ accessible, for $W$ to be Quillenizable on the whole of Cat. This is the case which seems to me at present the most interesting.

\textbf{Example 2:} Let $\mathcal{A}$ be an abelian category, and $\mathcal{M} = K^{\bullet}(A)$ the category of cochain complexes with values in $\mathcal{A}$, $W = W_{\mathcal{A}}$ the set of quasi-isomorphisms, hence % [sic] K^\bullet(A) -> K^\bullet(\mathcal{A}), plain italic A where calligraphic A is used everywhere else
\[ \mathcal{M}W^{-1} \;=\; K^{\bullet}(\mathcal{A})W_{\mathcal{A}}^{-1} \;\simeq\; D^{\bullet}(\mathcal{A}), \qquad \text{derived category of $\mathcal{A}$.} \]
We also define $\mathcal{M}^{+}$, $\mathcal{M}^{-}$, $\mathcal{M}^{\mathrm{b}}$, with the induced localizers $W_{\mathcal{A}}^{+}$, $W_{\mathcal{A}}^{-}$, $W_{\mathcal{A}}^{\mathrm{b}}$, and the categories of fractions $D^{+}(\mathcal{A})$, $D^{-}(\mathcal{A})$, $D^{\mathrm{b}}(\mathcal{A})$. Or
\msp{page 52}again, for $I$ in Cat,
\[ \mathcal{M}(I) \;=\; K^{\bullet}(\mathcal{A}(I)) \]
% foot: Version 25 janvier 2004
\printedpage{30}
% head: A. Grothendieck : Dérivateurs, Ch. XIII, Édition des Universités de Montpellier II et Paris VII
\[ W(I) \;=\; \underbrace{W_{\mathcal{A}(I)}}_{\text{quasi-isomorphisms}} \;\subset\; \mathrm{Fl}(K^{\bullet}(\mathcal{A}(I))), \]
and likewise for $\mathcal{M}^{\alpha}(I)$, $W^{\alpha}(I)$ $(\alpha \in \{+,-,\mathrm{b}\})$, at least \textit{if $I$ is finite.} There are also the subcategories $K^{\geq 0}(\mathcal{A})$ and $K^{\leq 0}(\mathcal{A})$ of cochain resp. chain complexes in the strict sense (with $K^{i} = 0$ if $i \geq 0$, resp. if $i \leq 0$ \texttt{[rather if $i < 0$, resp. if $i > 0$]}). We will have for arbitrary $I$
\[ \left\{ \begin{aligned} K^{\geq 0}(\mathcal{A})(I) &\;\simeq\; K^{\geq 0}(\mathcal{A}(I)) \\ W_{\mathcal{A}}^{\geq 0}(I) &\;=\; W_{\mathcal{A}(I)}^{\geq 0}, \end{aligned} \right. \]
and likewise for $K^{\leq 0}$, $W^{\leq 0}$.

According to \textsc{Quillen} page 1.2, $(K^{+}(\mathcal{A}),W_{\mathcal{A}}^{+})$ is Quillenizable if $\mathcal{A}$ has enough injectives, $(K^{-}(\mathcal{A}),W_{\mathcal{A}}^{-})$ is so if $\mathcal{A}$ has enough projectives. He says nothing for $(K^{\bullet}(\mathcal{A}),W_{\mathcal{A}})$ without more -- yet one does cohomological algebra with $D^{\bullet}(\mathcal{A})$ without more, and even without making an assumption of the existence of injectives or projectives. I have somewhat forgotten what \textsc{Verdier} did on this subject -- I should all the same look at his thesis (!), or else the exposition of \textsc{Hartshorne}.

\msp{page 53}At any rate, we have here examples of a quite different type from those considered in 1°. Thus, there is no need to assume that $\mathcal{M}$ is stable under small $\varinjlim$ -- $\mathcal{M}$ can even be a small category, it is so if and only if $\mathcal{A}$ itself is. If $\mathcal{A}$ is stable under (small) filtering $\varinjlim$, and is accessible (e.g. has a small generating family, and filtering $\varinjlim$ are exact), then $K^{\bullet}(\mathcal{A})$ etc. (\footnote[30]{!!}) have the same properties, and $(K^{+}(\mathcal{A}),W_{\mathcal{A}}^{+})$ is Quillenizable on the whole of Cat (\footnote[31]{careful, $K^{+}(\mathcal{A})$ is \textit{not} stable under infinite filtering $\varinjlim$ -- so one cannot apply the carpet of paragraphs 2, 3 directly to $K^{+}$.}). (This is the case where we know that $\mathcal{A}$ has enough injectives, and cohomological algebra is done mostly by dint of injective resolutions \dots) In this case we know that the % [sic] elle même -> elle-même, missing hyphen
\[ u_{*} \;:\; D^{+}(I',\mathcal{A}) \longrightarrow D^{+}(I,\mathcal{A}) \]
exist as right adjoints of the
\[ u^{*} \;:\; D^{+}(I,\mathcal{A}) \longrightarrow D^{+}(I',\mathcal{A}), \]
but I am no longer too sure whether we know how to define $u_{*}$ also on all of $D(I',\mathcal{A})$, and I doubt that $u_{!}$ exists without an assumption on $u$ (of the sort: finite cohomological dimension). Probably we must always have
\[ u_{!} \;:\; D^{-}(I',\mathcal{A}) \longrightarrow D^{-}(I,\mathcal{A}), \]
\msp{page 54}left adjoint to
\[ u^{*} \;:\; D^{-}(I,\mathcal{A}) \longrightarrow D^{-}(I',\mathcal{A}). \]
% foot: Version 25 janvier 2004

\printedpage{31}
% head: A. Grothendieck : Dérivateurs, Ch. XIII, Édition des Universités de Montpellier II et Paris VII
But I no longer really know whether we should expect to be able to define $u_{!}$ on all of $D(I',\mathcal{A})$, apart from assumptions like finite cohomological dimension. The sad truth is that I have $\pm$ forgotten the theory of derived functors! (I ought to ask \textsc{Jean Malgoire} again to send me canonical texts.)

\textbf{Scholie.} If $(\mathcal{M},W)$ is Quillenizable on Diag $\subset$ Cat (and even without its being Quillenizable), it defines a prederivator $\mathbf{D}_{\mathcal{M},W}$ on Diag. We will sometimes set
\[ \mathrm{HO}_{\mathcal{M},W}(I) \;\text{ or }\; \mathrm{HO}_{W}(I) \;=\; \mathbf{D}_{\mathcal{M},W}(I). \]
The Quillenizable hypothesis implies, on each of these homotopy categories, the (relatively weak) ``triangulated'' structure, coming from the \textit{exact} suspension and cosuspension sequences (\footnote[32]{In the case where $\mathcal{M}$ is pointed by \texttt{[a]} neutral object. Otherwise, there is the structure associated to cartesian and cocartesian squares. It would be quite interesting to axiomatize such a species of structure.}). (The Quillenizable character ensures the exactness of these sequences \dots) The functors $u^{*} : \mathrm{HO}_{W}(I) \longrightarrow \mathrm{HO}_{W}(I')$
\msp{page 55}are ``homotopically exact'' in the sense of these ``triangulated'' (noncommutative) structures, \textit{i.e.} transform exact suspension resp. cosuspension sequences into likewise. % [sic] des ces -> de ces, grammar

Which axioms of derivators are satisfied under the Quillenizable hypothesis alone? \textbf{Der} 1 \texttt{[see chapter I, p. 43]} is so without any hypothesis at all on the category of models $(\mathcal{M},W)$:
\[ \mathrm{HO}_{W}(\coprod_{\alpha} I_{\alpha}) \;\simeq\; \prod_{\alpha} \mathrm{HO}_{W}(I_{\alpha}) \]
for a \textit{finite} family of objects $I_{\alpha}$ in Diag. The localization axiom \textbf{Der} 2 \texttt{[see chapter I, p. 43]} is so under the hypothesis that $W$ be \textit{strongly saturated,} and Quillenizable. We use the fact that every arrow of $\mathrm{HO}_{W}(I)$ comes, up to isomorphism, from an arrow in $\mathcal{M}(I)$. (That was one of the points which remained pending in VI, in the theory of $\mathrm{HOT}_{W}$, $W \subset \mathrm{Fl}(\mathrm{Cat})$ a basic localizer \dots)

The axiom \textbf{Der} 3, new version \texttt{[see chapter I, pp. 111-112]}, involves a basic localizer $\mathbf{W} \subset \mathrm{Fl}(\mathrm{Cat})$ (not necessarily $\mathbf{W}_{\infty}$) (\footnote[33]{It is natural here to assume that $\mathbf{W}$ satisfies all the ``good axioms'' on basic localizers W (1, 2, 3, 4, 5, 6, 7) \texttt{[see chapter XII]}.}). The strongest form would be this: if
\[ u \;:\; I' \longrightarrow I \]
\msp{page 56}is in $\mathbf{W}$, we want the functor $u^{*}$ to induce an \textit{equivalence}
\[ u_{\mathrm{lc}}^{*} \;:\; \mathrm{HO}_{W}^{\mathrm{lc}}(I) \xrightarrow{\;\sim\;} \mathrm{HO}_{W}^{\mathrm{lc}}(I'). \]
But I doubt that we are equipped to verify so strong an axiom, without having at least a theory of the operations $u_{!}$ or $u_{*}$. One could on the off chance posit this axiom by taking $\mathbf{W} = \mathbf{W}_{\alpha} \overset{\text{def}}{=}$ the smallest basic localizer satisfying W (1 to 7) (the largest of
% foot: Version 25 janvier 2004
\printedpage{32}
% head: A. Grothendieck : Dérivateurs, Ch. XIII, Édition des Universités de Montpellier II et Paris VII
these localizers being $\mathbf{W}_{\infty}$). For the moment, the way to verify this axiom (in terms of the datum $(\mathcal{M},W)$) is for me a $\pm$ total mystery. (But I know how to verify it for $\mathrm{HOT}_{W}$, when $W \subset \mathrm{Fl}(\mathrm{Cat})$ satisfies W (1, 2, 3, 4, 5, \dots).)

I am going to look into it, once I have a reasonable command of the formalism of $u_{*}$, $u_{!}$ relative to the $u : I' \longrightarrow I$ in Diag. The Quillenizable hypothesis alone seems too coarse to ensure the existence of $u_{!}$ or of $u_{\star}$ in interesting cases. % [sic] u_\star -> u_*, star symbol where u_* is used elsewhere

\noindent\msp{page 57}

\section*{6\quad Quillenizable category of models, and functors $u_{!}$, $u_{*}$}

\textbf{Definition 1:}

\begin{itemize}[leftmargin=2.5em]
\item[\circled{a}] We say that $(\mathcal{M},W)$ is \textit{uniformly Quillenizable on the left} on Diag, if we can find a weak Quillen structure $(C,TC,F,TF)$ for $\mathcal{M}$ which is \textit{closed,} and which moreover has the following property: for $I$ in Cat, let $C(I)$, $TC(I) \subset \mathrm{Fl}(\mathcal{M}(I))$ be defined by
\[ \begin{aligned} &(u : X \longrightarrow Y) \in C(I) \text{ (resp. } TC(I)) \text{ if and only if } \ \forall\, s \in \operatorname{Ob} I, \\ &u_{s} : X_{s} \longrightarrow Y_{s} \text{ is in } C(I) \text{ (resp. } TC(I)). \end{aligned} \]
Let moreover
\[ F(I) \;=\; TC(I)_{*}, \qquad TF(I) \;=\; C(I)_{*}. \]
The condition in view is that for every $I$ \textit{in} Diag, $(C(I),TC(I),F(I),TF(I))$ is a \textit{closed exact} weak Quillen structure for $W(I)$ (\footnote[34]{$(C,TC,F,TF)$ is then called a \textit{uniform left Quillenization} of $(\mathcal{M},W)$ on Diag.}).

\item[\circled{b}] We say that $(\mathcal{M},W)$ is \textit{Quillenizable on Diag uniformly and exactly} \texttt{[on the]} \textit{left}, if we have in addition $TC = C \cap W$, \textit{i.e.} if $(C,TC,F,TF)$ is left exact, hence defined by $C$ alone, via
\[ TC = C \cap W, \qquad F = (TC)_{*}, \qquad TF = C_{*} \qquad \text{(\footnotemark[35])} \]
\footnotetext[35]{We say that $(C,\ldots)$, or simply $C$, defines a \textit{uniform and left exact Quillenization} of $(\mathcal{M},W)$ on Diag.}
\end{itemize}

We saw in paragraph 3 that if $\mathcal{M}$ is stable under small $\varinjlim$ and accessible, and $W$ stable under small filtering $\varinjlim$, finally if
% foot: Version 25 janvier 2004
\printedpage{33}
% head: A. Grothendieck : Dérivateurs, Ch. XIII, Édition des Universités de Montpellier II et Paris VII
\msp{page 58}$W$ admits a weak Quillen structure $(C,TC,F,TF)$ which is closed and left exact -- then the conditions of definition 1 are satisfied.

On the other hand, let us note that proposition 1 page 45 and its corollaries, including proposition 3 (page 50) and its corollaries, are valid provided that the Quillenizations of the $(\mathcal{M}(I),W(I))$ under consideration are defined ``uniformly on the left'' by a uniform left Quillenization (not necessarily left exact) of $(\mathcal{M},W)$ on Diag.

Uniform left Quillenization will serve us to establish the existence of homotopy direct images $u_{*}$, for certain $u$. Uniform right Quillenization, dually, will have as its main consequence the existence of homotopy $u_{!}$ for certain $u : I' \longrightarrow I$ in Diag.

I am going to summarize what I have been able to establish so far in this direction, in the

\msp{page 59}\textbf{Theorem-scholie:} \textit{Let $(C,TC,F,TF)$ be a uniform left Quillenization of $W$ on} Diag $\subset$ Cat.

\begin{itemize}[leftmargin=3em]
\item[1)] \textit{Let $u : I' \longrightarrow I$ be an arrow in} Diag. \textit{We assume} Diag \textit{matched to $\mathcal{M}$ in such a way that for $u$, an arrow in} Diag, \textit{the $\varprojlim$ in $\mathcal{M}$ (of type $I'/i$) required to ensure the existence of}
\[ u_{*} \;:\; \mathcal{M}(I') \longrightarrow \mathcal{M}(I) \]
\textit{exist in $\mathcal{M}$. (This is O.K. if $I$, $I'$} \texttt{[are]} \textit{finite, $\mathcal{M}$ being assumed stable under finite $\varinjlim$ in any case.) We know} (prop. 1, page 45) \textit{that $u_{*}$ induces}
\[ u_{*}^{\mathrm{f}} \;:\; \mathcal{M}(I')_{\mathrm{f}} \longrightarrow \mathcal{M}(I)_{\mathrm{f}}\,. \tag*{$(*)$} \]
\textit{We say that $u$ is} admissible (on the left) for the given uniform left Quillenization of $W$ on Diag, \textit{if $u_{*}^{\mathrm{f}}$ is compatible with the localizers $W(I')_{\mathrm{f}}$, $W(I)_{\mathrm{f}}$,} i.e. \textit{if $f' : X' \longrightarrow Y'$ in $\mathcal{M}(I')$, $X'$, $Y'$ fibrant, $f' \in W(I')$ implies $u_{*}(f') \in W(I)$. We say that $Z$ in} Diag \textit{is} admissible (for the given uniform left Quillenization of $W$ on Diag \dots) \textit{if the functor $Z \longrightarrow e$ is admissible. This being so:}
\begin{itemize}[leftmargin=2.5em]
\item[a)] \textit{For $u$ to be admissible, it suffices that the $I'/s$ for $s \in \operatorname{Ob} I$, be so. NB there is a converse, when $I$ is ordered, hence the $I'/s \longrightarrow I'$ open immersions.}
\item[b)] \msp{page 60}\textit{Every finite ordered set is admissible. More generally (barring error) every ordered set $I$ such that the topological space associated to $I^{\circ}$ is noetherian (}i.e. \textit{every increasing sequence of open sets of $I^{\circ}$ is stationary, or again every increasing sequence of closed sets of $I$ is stationary), and if moreover every irreducible closed subset of $I^{\circ}$ admits a generic point,} i.e. \textit{$I^{\circ}$ is (not only a noetherian space, but also) a topological space which is} sober. \textit{(For $I$, this means that every upward filtering open subset of $I$ admits a greatest element.)} % [sic] fermées -> fermés, agreement with masculine (ouverts / fermés)
\end{itemize}
\end{itemize}
% foot: Version 25 janvier 2004